\documentclass[12pt]{article}
\usepackage[margin=1in]{geometry}
\usepackage{enumitem}
\usepackage{titlesec}
\usepackage{parskip}
\usepackage{graphicx}
\usepackage{xcolor}
\usepackage{hyperref}
\usepackage{fontenc}

\title{\textbf{Suryanarayan 2024 RG Notes}\\
\large Suryanarayan, Pavithra. ``Endogenous State Capacity'' [2024].\\
\textit{Annual Review of Political Science}.}
\author{}
\date{}

\begin{document}

\maketitle

Three primary assumptions of state capacity that have been historically held -- recently, some (but not total) shift away from them:

``The first is that state capacity is inordinately complex. To have strong capacity requires several types of organizational and technical competencies firing at once---information, administration, taxation, policing, and legal prowess. The second is that capacity is hard to build, in part because of the complexity of the endeavor, but also because rulers in the process of building capacity encounter resistance from society---the elites and masses. The third assumption is that once a stock of state capacity emerges, it is relatively stable, and the multiple dimensions of capacity---coercive, extractive, and bureaucratic---usually move in tandem over time.'' (Suryanarayan, p. 224)

Movement towards endogenous mechanism, noting that the involved actors have their own rational incentives for various amounts of state capacity

New opportunities for research (some existing, but still small) in impact of identity and culture on state capacity formation

\end{document}
